% Page 004 — Carte du pays de Meyrueis
% Page liminaire : image pleine page, sans pagination.

\thispagestyle{empty}
\sectiondate{À propos de cette édition}{}

Lorsque nous avons découvert, au détour d'une recherche sur le village de Meyrueis, ce recueil nous l'avons lu avec beaucoup d'attention.

Préface.

Note de l'éditeur — ou « À propos de cette édition ».

Guide de lecture et conventions typographiques.

Liste des abréviations.

Table des matières.

Texte du livre.


1. Préface
Le sens et l'esprit du livre
- Pourquoi ce livre mérite d'être relu aujourd'hui.
- La mémoire collective de Meyrueis et de ses habitants.
- La valeur des témoignages, des documents et des souvenirs locaux.
- La place de Paul Loupiac dans cette transmission.
- Ce que le lecteur peut en retirer, qu'il connaisse ou non le village.




2. À propos de cette édition
Expliquer ton travail éditorial
- Origine du document utilisé et nature de la reproduction.
- Transcription et remise en page.
- Respect de l'orthographe, du vocabulaire et des formulations d'origine.
- Traitement des erreurs manifestes, des passages incertains et des éventuelles lacunes.
- Ajouts éditoriaux : notes, légendes, liens, illustrations ou informations complémentaires.




3. Guide de lecture
Aider à comprendre la structure
- Organisation en sections, sous-sections et sous-sous-sections.
- Fonction de la table des matières et des dates associées aux titres.
- Présentation des notes de bas de page.
- Particularités des documents anciens reproduits dans le corps du livre.




4. Abréviations, conventions et typographie
Un aide-mémoire concis
- Abréviations réellement employées dans le livre.
- Signification des sigles et des termes historiques ou locaux.
- Conventions pour les dates, les siècles et les périodes.
- Signes utilisés pour signaler une précision éditoriale, une incertitude ou un passage manquant.




5. Repères géographiques et historiques (facultatif)
- Carte de situation de Meyrueis et des vallées évoquées.
- Noms anciens et noms actuels des lieux, si nécessaire.
- Petit repère chronologique pour les périodes qui reviennent souvent.

\subsection{Abréviations utilisées dans ce recueil}

\begin{itemize}
    \item\textbf{Dr.} Docteur
    \item\textbf{F} Franc(s)
    \item\textbf{Mme} Madame
    \item\textbf{Mgr} Monseigneur
    \item\textbf{M.} Monsieur
    \item\textbf{Sr.} Sieur
\end{itemize}
